% =================================================================
\section{Dataset \& Problem Analysis}
\label{sec:data}
% =================================================================

\subsection{Dataset}
\label{sec:data:dataset}

The CD-MSC corpus contains 271,380 audio clips of nine mosquito
species recorded across five conditions (D1--D5).
The labels D1--D5 are opaque---the challenge release does not
include device or location metadata---so we treat them as five distinct
but uncharacterised recording environments~\cite{biodcase2026}.
All clips are resampled to 8~kHz (the challenge-specified rate) and
converted to 64-bin log-mel spectrograms using the official baseline
frontend: 512-sample Hann window (25~ms), 80-sample hop (10~ms),
mel filterbank from 0 to 4{,}000~Hz.
Because mel frequency spacing is non-linear, the bin-to-frequency
mapping ranges from $\sim$36~Hz per bin at the low end to
$\sim$140~Hz per bin at the high end.
Features are normalised by subtracting the per-bin mean and dividing
by the per-bin standard deviation, both computed from the training
split only.
Clip lengths are mostly 0.63~s (mode = 63 frames); the test set
contains some multi-second recordings.
The official training/validation/test split is 213,647 / 30,516 /
27,217 clips.

\subsection{The Core Problem: One Domain Dominates Everything}
\label{sec:data:imbalance}

Table~\ref{tab:domain_counts} and Fig.~\ref{fig:data_analysis}~(left)
show the training-set breakdown.
D5 holds \textbf{212,339 of the 213,647 training clips---99.4\%}.
Conditions D1--D4 together contribute only 1,308 clips,
split across 80 to 634 clips each.

\begin{table}[t]
  \centering
  \caption{Training-set domain distribution. D5 holds 99.4\% of
           training data; D1--D4 together account for only 1,308 clips.}
  \label{tab:domain_counts}
  \setlength{\tabcolsep}{5pt}
  \begin{tabular}{lrrrr}
    \toprule
    Condition & Train & Val & Test & Train \% \\
    \midrule
    D1 & 634   & 96     & 3{,}335  & 0.30 \\
    D2 & 230   & 30     & 524      & 0.11 \\
    D3 & 364   & 53     & 262      & 0.17 \\
    D4 & 80    & 3      & 117      & 0.04 \\
    D5 & 212{,}339 & 30{,}334 & 22{,}979 & \textbf{99.38} \\
    \midrule
    Total & 213{,}647 & 30{,}516 & 27{,}217 & 100 \\
    \bottomrule
  \end{tabular}
\end{table}

\paragraph{Why this breaks domain-generalisation methods.}
Domain-adversarial training (DANN) works by training a discriminator
to classify which domain a sample came from, then reversing its
gradients to push the model toward domain-agnostic features.
But when 99.4\% of batches come from D5, the discriminator's easiest
strategy is to learn one thing: ``is this D5 or not?''
Reversing that gradient removes only the D5-vs-not-D5 signal, leaving
all D1--D4-specific features in the embedding intact.
When training data is balanced across domains, the discriminator
must separate all five conditions, and the reversed gradient removes
all of them.
This mechanism explains our central empirical finding: DANN without
balanced training gains only +0.4~pp $\BAu$; DANN \emph{with}
balanced training gains +20.3~pp.

\paragraph{Why D5 is excluded from our evaluation.}
Holding D5 out in a Leave-One-Domain-Out experiment would mean
training on only 1,308 clips---roughly the size of a small pilot
study---which is a data-starvation problem, not a domain-shift problem.
We exclude D5 from all LODO means; all reported numbers are D1--D4
mean.

\subsection{Why the Official Evaluation Metric Is Misleading}
\label{sec:data:split}

We compute the domain centroid of each condition in the training set
(the per-domain average log-mel feature vector) and assign each test
clip to its nearest centroid by Euclidean distance.
Approximately 85\% of test clips are closest to the D5 centroid.
The official 10-seed challenge baseline reports
$\BAu{=}0.175$~\cite{biodcase2026}---but 85\% of those ``unseen''
test clips effectively look like D5, making this a measure of
\emph{D5 generalisation}, not cross-condition transfer.
Our LODO baseline (seed~42, D1--D4 mean) gives
$\BAu{=}0.165$, $\BAs{=}0.578$, $\DSG{=}0.412$.
\textbf{We use LODO D1--D4 mean $\BAu$ as our primary metric
throughout.}

\subsection{Which Frequency Bands Carry Species vs.\ Recording-Condition Information?}
\label{sec:data:freq_analysis}

To decide \emph{where} in the spectrogram to apply data augmentation,
we ask: which mel bins vary more \emph{within} one recording condition
(species-specific variation) than \emph{across} conditions
(condition-specific variation)?
We compute this ratio---within-domain variance divided by
cross-domain variance---for each of the 64 mel bins, using the
held-out test split.\footnote{This analysis uses test-split data solely
to motivate the design choices $K{=}9$ and $K'{=}36$.
Both values were fixed before any FBS-Mix experiments; no further
tuning was performed.}

A ratio below~1 means the bin varies more across recording conditions
than within them: it is picking up \emph{recording noise}, not
species signal.
A ratio above~1 means it varies more within a condition: it is
picking up something that differs between species, not between
recording devices.

Fig.~\ref{fig:data_analysis}~(right) shows three clear zones:

\begin{itemize}
  \item \textbf{Bins 0--8} ($\approx$36--325~Hz), ratio 0.30--0.93:
    condition-dominated.
    These bins capture low-frequency recording noise and microphone
    response---information that changes with the device, not the mosquito.
  \item \textbf{Bins 9--35} ($\approx$361--1{,}360~Hz), ratio 1.50--4.12:
    species-dominated.
    This is the mosquito wingbeat core: the ratio peaks at bin~17
    ($\approx$649~Hz), near the wingbeat fundamental of the most
    common species~\cite{mukundarajan2017}.
  \item \textbf{Bins 36--63} ($\approx$1{,}412--3{,}854~Hz),
    ratio 0.73--1.47: mixed---upper harmonics and background noise
    are tangled together.
\end{itemize}

\textbf{Why mosquito wingbeat is device-invariant.}
Mosquito wings beat at a species-specific rate (300--800~Hz
fundamental, with harmonics extending to $\sim$1{,}400~Hz)
determined by wing morphology~\cite{mukundarajan2017}.
This physical oscillation is independent of the recording device, so
the wingbeat frequency content stays consistent across conditions even
when everything else changes.

This analysis tells us precisely what MixStyle gets wrong.
MixStyle~\cite{mixstyle2021} mixes the statistics of \emph{all}
frequency channels inside the CNN, which randomises the wingbeat
core (bins 9--35) at the same time as the condition-noise band
(bins 0--8).
Corrupting species information cancels any benefit from diversifying
condition noise---which is why balanced+MixStyle and balanced-only
perform identically ($\approx$0.32~$\BAu$).
FBS-Mix~(Section~\ref{sec:method:fbsmix}) mixes only the
condition-dominated band and leaves the wingbeat core alone.
